\section{Which source could this actor use at this date?}
\label{sim:source-timed-policies}

The four feasible policies in the preceding example depend on two supplied interpretations: which receipt and facility-right family is possible, and which signals the decision maker observes. A revised document can change either interpretation. It can also leave both interpretations unchanged while correcting the evidence supporting them. The simulator needs to distinguish those cases before it decides which calculation to repeat.

Consider a correction released on 3 January to a document used for a decision on 2 January. The correction may improve a retrospective account of the company, but it cannot be inserted into what the actor knew on 2 January merely because the simulator downloads it later. Conversely, a document may have been available to the actor on 1 January but entered the simulator on 3 January and passed internal review on 4 January. On 4 January the system can reconstruct the earlier information set if the retained evidence supports that access. It has not established that the reconstruction itself was run on 2 January.

These are distinct dates with distinct subjects. Publication concerns the document's release. Access concerns a particular actor. Ingestion and review concern the information system. Effective time concerns the economic period or event to which the interpretation applies. A publication timestamp alone does not establish personal access; a download timestamp does not establish either the original publication or the document's economic effect.

\subsection{Select revisions before selecting the policy problem}

The selected construction uses an immutable record for each supplied interpretation. The record identifies the source bytes, the interpretation's subject, the previous revision it replaces, its effective time, and separate observations of publication, access, ingestion and review. The caller supplies the evidence accepted for these timing observations. Hashes preserve the identity of those bytes and observations; they do not establish their truth.

For a request, retain revisions that are effective for the chosen economic event, available to the actor by the decision time, and ingested and reviewed by the system's reconstruction cutoff. Known later revisions are excluded from that earlier view. A missing relevant time leaves the selection unresolved. Within a subject, select a revision only if it is the unique remaining head of its replacement chain. Two incomparable available corrections preserve a conflict. Retraction of a required interpretation leaves a missing input rather than silently restoring its earlier assertion.

This procedure is deliberately restricted to a declared source scope. An unrelated management note cannot change the coupon interpretation merely because it appears in the same registry. A relevant correction must be inspected again on each request, however. Comparing a cached view identifier without checking for newly available revisions could miss exactly the change that should invalidate the result.

In the reference example, one required subject supplies receipt 6 or 7 and the mixed or committed facility-right family; the other supplies the coarse or refined observation scheme. These choices map to the eight previously certified problems. They preserve the same possible persistent contract types, complete financial rows and rejected actions within the selected family. No convenient prior or likelihood is introduced to make a source selector callable. The question remains whether one policy works in every history the actor cannot distinguish.

The timing checks reuse the existing source and policy-information mechanisms. The selected adapter accepts explicit whole-second UTC instants and public-document evidence with a separate supplied actor-access observation. It refuses finer or incomplete clock representations instead of rounding them into eligibility. Private communications and more general time intervals need their own extension. An unresolved access record is therefore a substantive result about this request, not an instruction to substitute the publication date.

\subsection{Reuse a calculation without reusing an obsolete source binding}

Suppose a replacement document supports exactly the same receipt, facility-right family and observation scheme. The set of robust policies can remain unchanged. Nevertheless, a report that still cites the replaced source would give the wrong explanation of which evidence supported the current request. Numerical equality and evidential identity answer different questions.

Let $F$ denote the retained financial relations and certified rows, $s$ the selected receipt and facility family, and $o$ the observation scheme. Write the complete reachable policy calculation as
\begin{equation}
 \mathcal P(F,s,o).
 \label{eq:source-policy-dependency}
\end{equation}
Its numerical inputs do not include a document's filename or the system's ingestion date. Those records determine whether the interpretation $(s,o)$ is available for the actor and request. The reported result therefore needs both $\mathcal P(F,s,o)$ and the source, revision, actor and cutoff records that justify using these inputs.

The checked dependency relation makes this separation explicit. Certified financial evidence supports the chosen financial scenario; that scenario and the observation scheme support the policy set. Separate source bindings attach the selected financial and observation interpretations to their revision histories. The final actor binding joins the numerical policy set to both source bindings and to the decision, reconstruction and effective cutoffs.

A same-meaning correction changes a source binding and the final actor binding while preserving the numerical policy set. A change from coarse to refined information preserves the financial relations but requires a new observation and policy calculation. A changed receipt or facility family changes the financial scenario and its dependent policy problem. Changing only the actor or cutoff requires a new actor binding even when the selected interpretations and numerical answer happen to be identical. A correction known to arrive after the old cutoff changes none of that earlier view's selected dependencies.

Reuse is justified by those dependencies, not by observing that the new and old policy counts coincide. The reference compares the entire reused policy set and its complete world-row population with a fresh calculation. It also requires the saved inputs and method bytes to match expected identities supplied outside the transported bundle. Removed or duplicated rows, changed hidden rights or signals, altered policy references and rewritten prior results are rejected. Recomputing an inner checksum after changing a saved result does not make it the previously accepted result.

The scope remains conditional. The caller must establish which source and timing evidence is accepted, and a complete history family for an actual company may require many more subjects than these two. The result shows how revision-sensitive selection can coexist with numerical reuse while preserving the robust question. It establishes neither authenticated storage nor crash recovery, and the small exact comparison supplies no performance ranking.\footnote{The immutable source-view cases are retained in \path{data/accounting-master/certified-source-views-20261026/}; the separate dependency and complete cold-versus-reused comparisons are in \path{data/accounting-master/source-policy-invalidation-20261027/}. Both preserve their supplied source scope and withhold actual-source, accounting, real-actor and forecast admission.}
